\section{Background}

Denoising diffusion probabilistic models learn to generate data by reversing a fixed forward
process that progressively corrupts clean samples with Gaussian noise. Formally, the model
trains a neural network $\epsilon_\theta(x_t, t)$ to predict the noise injected at diffusion
step $t$, where $x_t$ is obtained by sampling from the forward noising process applied to a
clean input $x_0$ \cite{ho2020ddpm}. At inference time, new samples are generated by
iteratively applying the learned denoising network from a pure Gaussian prior, guided by a
fixed noise schedule. The quality of this reverse process depends critically on the capacity
of $\epsilon_\theta$ to model the score function of the data distribution across all noise
levels---a requirement that motivates the use of expressive U-Net and transformer-based
backbone architectures \cite{ho2020ddpm, peebles2023dit}.

LoRA adapts a pretrained weight matrix $W \in \mathbb{R}^{d_{\text{out}} \times d_{\text{in}}}$
by introducing a low-rank residual update of the form
\begin{equation}
  W' = W + BA,
  \label{eq:lora}
\end{equation}
where $A \in \mathbb{R}^{r \times d_{\text{in}}}$ and $B \in \mathbb{R}^{d_{\text{out}} \times r}$
are the trainable adapter matrices, and $r \ll \min(d_{\text{out}}, d_{\text{in}})$ is the
rank \cite{hu2022lora}. During fine-tuning, the base parameters $W$ are frozen and only $A$
and $B$ are optimized, substantially reducing the number of trainable parameters relative to
full fine-tuning. The total adapter parameter count scales linearly with $r$: increasing rank
expands the expressivity of the low-rank update but proportionally increases trainable parameter
counts, memory footprint, and, to a lesser extent, runtime. Crucially, the adapted weight
$W'$ takes the same form as the original matrix, so trained adapters can be merged into $W$
at inference with no additional computational overhead.

Because rank directly governs trainable capacity, comparing LoRA configurations across
different ranks is a non-trivial experimental design problem. If optimizer settings---such as
learning rate, batch size, or training duration---are adjusted separately for each rank, the
observed quality differences conflate the effect of rank with the effect of tuning. Our study
therefore holds all optimization hyperparameters fixed across ranks, interpreting observed FID
differences as evidence of rank efficiency under a controlled and uniform training budget,
rather than as indicators of per-rank hyperparameter optima. This design choice is central to
the empirical framing of all results reported in subsequent sections.